\chapter{Against Analysis}
Madoka Magica the Movie 3: Rebellion is a difficult subject to approach. Like the series, it is semiotically dense (that is, the images and events it contains evoke a relatively high number of readings); however, where the series primarily takes place in fairly stark, sparse, almost sterile environments, to contrast the visual complexity of the witch's labyrinths, the movie is tremendously visually complex throughout. In other words, not only does each image have as much to unpack as any image in the series, it's also got quite a bit more imagery than its five-and-a-half episodes' worth of runtime would suggest!

Analyzing it is thus a potentially overwhelming task. To take the movie as an organic whole is nigh-impossible, at least in anything like chapter length; to do it justice, one must either pick a theme and follow its development through the film, missing out on all the other themes except perhaps for how they interact with the chosen theme, or one can pick a scene and examine it in all its complexities, missing out on all other scenes except insofar as they impact the chosen scene.

Consider a relatively straightforward question: What is the titular rebellion? Is it Kyubey against Madoka's new world? Homura against the world inside the labyrinth, unaware that she herself created it? Is it Homura against Madoka? Or is the film itself an act of rebellion, and if so, who or what is it rebelling against?

Even a question as seemingly simple as "What does the title mean?" leads only to more questions, and there are a multitude of supportable answers to each of them, each of which could fill an essay in its own right. Thus, the remaining chapters of this book comprise a series of essays on Rebellion. Some may be analyses of single scenes; others may trace themes or look at the evolution of a character. All, however, are explorations of a particular answer to the question, "Against what?" Such an approach seems to me the only viable way in which I can approach Rebellion, as this is a film that defies analysis.

And a near-synonym for defiance is... well, you get the idea.

So, then, let us consider a particular short scene that exemplifies how difficult this film is to analyze, specifically, the final post-credits stinger. Immediately prior to the stinger, the credits themselves depict a heavily stylized version of the movie's plot, with Homura and Madoka divided by the credit scroll itself. At the very end of the credits, however, they hold hands and run off together into the distance, a surprisingly hopeful end to the story given Homura's posturing in the final scenes before the credits. Is this foreshadowing, or just Homura's dream? Is the fact that they vanish into the distance evidence that they will escape, or evidence that the possibility of them being together is disappearing?

It doesn't matter, because the stinger contradicts the image anyway. (Or does it? If the credits are Homura's dream and the stinger the reality, or the credits are foreshadowing and the stinger her fears...) It's worth, here, examining the usual function of a stinger. Most commonly found in big blockbuster action franchises (the Marvel Cinematic Universe has raised them to an art form) or comedies, the usual function of a stinger is either to serve as a punchline to a joke set up earlier in the film (possibly the best example of this is the taxi passenger in Airplane! \cite{ref113}) or to build excitement for and drop hints regarding the plot of the next installment in the franchise (Samuel L. Jackson would like to talk to you about the Avengers Initiative \cite{ref114}). 

Here, however, the stinger's function appears to be neither. Rather, its main function appears to be to generate questions and cast doubt on the way Homura chose to present herself in the final scenes prior to the credits. Admittedly, it is a "punchline" in the sense of concluding a repeated motif throughout the Madoka Magica movies, specifically the two chairs in a field. In the opening credits of the two compilation movies, we see Madoka and Homura sitting side-by-side on white chairs in the middle of a field of grass and flowers, cuddling playfully. In the Rebellion stinger, Homura begins sitting in a similar chair, alone, positioned on the edge of a cliff. Chair, cliff, and half-moon are lined up to create the effect of a picture sliced in two, as if the other half of the moon and the other half of the world, including Madoka and her chair, have been simply cut away and replaced with empty darkness.

Is this Homura's decision to make even Madoka her enemy? Her regret? Or just a cruel reminder for the audience of what has been lost?

Homura, in the final scenes of the movie, appears to be in near-total control. An army of familiars obey her; she can rewrite Sayaka's memories and cut her off from her Oktavia form; she can block Madoka from her Buddha-nature, the Law of Cycles. She is the creator of this new world, having rewritten reality earlier in the movie; it is not too far-fetched to suggest that she is now the most powerful entity within the confines of the universe (it is up for grabs how she compares to the Law of Cycles, as the two entities she appears unable to completely control are herself and Madoka—note her apparent panic when Madoka begins remembering her Buddha-nature shortly before the ending credits).

Why, then, does she appear startled by the approach of Kyubey in the stinger? The expression on her face is readable as either apprehension or hope; given the associations of the chair, does she momentarily believe it's Madoka? Does she hope it is, or fear that it is? How can she not know that it's Kyubey?

And then there is Kyubey's state: disheveled, trembling. Extreme close-ups on Kyubey's eye were frequently used in the series and this film to remind the viewer that he is watching, and they thus served to make him a more ominous and menacing figure. This close-up, however, shows his fur matted, his eye dulled, darkened, and shaking. For the first time, he displays a change in affect, but this serves only to unsettle the viewer further—even Kyubey, in this scene, is transformed into something "wrong." He is no longer a menacing figure, but rather a pathetic one, beaten and broken by Homura's display of power in rewriting the universe. This is a worst-case scenario for him and his kind; Madoka came to fear and distrust him, but she has little capacity for hate. Homura is different; full of rage and sorrow, it would not be at all out of character for her to take that out on the Incubators in general and the instance of Kyubey in Mitakihara in specific.

But the logic of the stinger suggests that the extreme close-up of Kyubey's eye is foreshadowing—it is the most typically stinger-like of any shot in this sequence, reminiscent of a horror movie ending with the believed-dead killer's eyes snapping open. Unfortunately, he is as inscrutable as ever; is he plotting a counterstroke against Homura? Simply observing and biding his time? Or is he truly broken, his pathetic appearance evidence that his role as villain has been stripped from him by Homura?

And then there is the dance. Homura dances with her new Soul Gem, both the style and music reminiscent of her balletic transformation sequence near the beginning of the movie. The gem resembles the chess symbol for a king, a near-powerless piece which is nonetheless the most important in the game. Is this Homura as queen—the strongest piece on the board—dancing with the most important figure in her world, the now-impotent Madoka? This reading is supported by the fact that the Soul Gem in question was made earlier in the movie from the pieces of Homura's old Soul Gem and a spool of thread the same color as Madoka's hair, but only if we read that thread as signifying Madoka herself or her connection to Homura. An equally likely reading is that it represents Madoka-the-incarnate-person's connection to Madoka-the-omnipresent-intangible-abstraction, in which case Homura is not so much missing her "other half" as reveling in her imprisonment. The difference between the two readings, in other words, lies between reading Homura as putting on a bold face over confusion and pain, or as a creepy, controlling stalker.

At the end of her dance, Homura tips sideways over the cliff. Her pose as she falls recalls Madoka's similar sideways tip off her chair when Homura becomes a witch, which seems fairly clearly to be a reference to Madoka's self-sacrifice and Homura's growing regret at failing to stop her. So is Homura seeking to join Madoka by replicating her action? Sacrificing herself so that Madoka doesn't have to? Mocking Madoka's sacrifice as a signifier that she has descended so far into evil that even the love which motivated her descent no longer matters? Or is it a futile gesture toward an impossible suicide (it is unclear what would happen to the universe if Homura died, but virtually certain that at least Madoka would reconnect with the Buddha-nature Homura is determined to keep her from) by a character in the depths of despair? Or is this (as well as or instead of other readings) a reference to the Tarot card The Fool, which frequently depicts the titular Fool about to step off a cliff, accompanied by a small white animal \cite{ref115}? 

We could explore these questions in detail, certainly, along with other questions (for example, the significance of the moon being precisely halfway between the almost-new moon when Sayaka became a witch and the full moon when Homura became one). It would take thousands of words and produce no certain conclusion except that the scene is deliberately ambiguous, but it can be done. That's not the point. The point is that this is ninety seconds of a two-hour movie, and not even the visually or semiotically densest ninety seconds (those, I suspect, fall somewhere between Homura witching out and she and Madoka shattering the Incubators' barrier). 

No, the point is this: This movie is dense, and it is ambiguous, and it thus poses a challenge to analysis.

Good. Let's do it anyway.
